% ======================================================
%  CHEETAH Program Sheet

%  SECTION TO SOURCE MAP
%    The program, layer by layer  systems/cheetah_joint/joint_lp.py row
% ======================================================

\documentclass[10pt]{article}

\newif\ifPLEX \PLEXfalse   % true + lualatex if IBM Plex is installed

\usepackage{iftex}
\ifPLEX
  \usepackage{fontspec}
  \setmainfont{IBM Plex Serif}
  \setsansfont{IBM Plex Sans}
  \setmonofont{IBM Plex Mono}[Scale=MatchLowercase]
\else
  \usepackage[T1]{fontenc}
  \usepackage{lmodern}                  % keeps Latin Modern Mono as \ttfamily
  \usepackage[scale=0.96]{tgheros}     % sans, body text
  \usepackage{tgschola}                 % serif, headings
\fi
\renewcommand{\familydefault}{\sfdefault}

\usepackage[a4paper,left=16mm,right=16mm,top=15mm,bottom=17mm]{geometry}
\usepackage[dvipsnames,table]{xcolor}
\usepackage{amsmath}
\usepackage{array}
\usepackage{booktabs}
\usepackage{adjustbox}
\usepackage{enumitem}
\usepackage{titlesec}
\usepackage{parskip}
\usepackage{microtype}
\usepackage{tikz}
\usepackage{tcolorbox}
\tcbuselibrary{skins,breakable,raster}
\usepackage{needspace}
\usepackage{xurl}
\usepackage[hidelinks]{hyperref}
\usepackage{bookmark}

% ---------------------------------------------------------------- dates
\newcommand{\sheetdata}{2026-10-01}
\newcommand{\sheetrev}{2026-10-02}

% --------------------------------------------------------------- palette
% Light theme of the artifact. The artifact also carries a dark theme;
% a PDF has one background, so only the light values are kept here.
\definecolor{paper}{HTML}{F5F7F9}
\definecolor{surface}{HTML}{FFFFFF}
\definecolor{surf2}{HTML}{EEF2F6}
\definecolor{ink}{HTML}{121B26}
\definecolor{ink2}{HTML}{465566}
\definecolor{ink3}{HTML}{6D7D8E}
\definecolor{rulec}{HTML}{CED6E0}
\definecolor{rulesoft}{HTML}{E2E8EE}
\definecolor{accent}{HTML}{1D5C8F}
\definecolor{accentsoft}{HTML}{E9F0F7}
\definecolor{sone}{HTML}{2A78D6}
\definecolor{stwo}{HTML}{EB6834}
\definecolor{okc}{HTML}{1F6F4A}
\definecolor{oksoft}{HTML}{E4F1EA}
\definecolor{warnc}{HTML}{8A6213}
\definecolor{warnsoft}{HTML}{F7EFDC}
\definecolor{stopc}{HTML}{9E3127}
\definecolor{stopsoft}{HTML}{F7E6E3}

\pagecolor{paper}
\color{ink}
\hypersetup{colorlinks=true,urlcolor=accent,linkcolor=accent,
            pdftitle={CHEETAH Program Sheet},
            pdfsubject={Hardware and software co-design, program reference}}

% ------------------------------------------------------------- utilities
\newcommand{\mono}{\ttfamily}
\newcommand{\tx}{\,\ensuremath{\times}}
\newcommand{\mmsq}[1]{#1\,mm\textsuperscript{2}}
\newcommand{\eyebrow}[1]{{\mono\fontsize{7.6}{9}\selectfont\color{ink3}\MakeUppercase{#1}}}
\newcommand{\lede}[1]{{\color{ink2}#1\par}\vspace{2pt}}

% inline identifier box, underscores may be raw
\newtcbox{\idbox}{on line,arc=1pt,outer arc=1pt,boxrule=0.3pt,boxsep=0pt,
  left=1.7pt,right=1.7pt,top=0.7pt,bottom=0.7pt,
  colback=surf2,colframe=rulesoft}
\newcommand{\idt}[1]{\idbox{\mono\fontsize{7.8}{9}\selectfont\detokenize{#1}}}

% row family tag
\newtcbox{\fambox}{on line,arc=1pt,outer arc=1pt,boxrule=0.3pt,boxsep=0pt,
  left=1.5pt,right=1.5pt,top=0.6pt,bottom=0.6pt,
  colback=surf2,colframe=rulesoft}
\newcommand{\rowfam}[1]{\fambox{\mono\fontsize{6.8}{8}\selectfont\color{ink2}\detokenize{#1}}\hspace{1pt}}

% provenance chips
\newtcbox{\chipbox}[1]{on line,arc=1.2pt,outer arc=1.2pt,boxrule=0.3pt,boxsep=0pt,
  left=2.4pt,right=2.4pt,top=1pt,bottom=1pt,#1}
% \texorpdfstring keeps the chip out of the PDF bookmark strings, since a
% chip may appear inside a \section title.
\newcommand{\chipgen}[3]{\texorpdfstring{%
  \chipbox{colback=#2,colframe=#1,coltext=#1}%
  {\mono\fontsize{6.4}{7}\selectfont\MakeUppercase{#3}}}{}}
\newcommand{\chipM}{\chipgen{okc}{oksoft}{measured}}
\newcommand{\chipR}{\chipgen{accent}{accentsoft}{report}}
\newcommand{\chipN}{\chipgen{warnc}{warnsoft}{notes}}
\newcommand{\chipL}{\chipgen{stopc}{stopsoft}{lost}}

% verdict badges
\newcommand{\vbadge}[3]{\chipbox{colback=#2,colframe=#2,coltext=#1}%
  {\mono\fontsize{6.6}{7}\selectfont\MakeUppercase{#3}}\hspace{3pt}}
\newcommand{\vyes}[1]{\vbadge{okc}{oksoft}{#1}}
\newcommand{\vno}[1]{\vbadge{stopc}{stopsoft}{#1}}
\newcommand{\vpart}[1]{\vbadge{warnc}{warnsoft}{#1}}

% ------------------------------------------------------------- headings
\setcounter{secnumdepth}{0}
\titleformat{\section}{\rmfamily\bfseries\fontsize{15.5}{18}\selectfont\color{ink}}{}{0pt}{}
\titlespacing*{\section}{0pt}{16pt}{3pt}
\titleformat{\subsection}{\sffamily\bfseries\fontsize{10.5}{13}\selectfont\color{ink}}{}{0pt}{}
\titlespacing*{\subsection}{0pt}{11pt}{2pt}

% ---------------------------------------------------------------- boxes
\tcbset{cardstyle/.style={enhanced,breakable,colback=surface,colframe=rulec,
  boxrule=0.4pt,arc=2pt,left=8pt,right=8pt,top=6pt,bottom=6pt}}
\newtcolorbox{card}{cardstyle}
\newtcolorbox{vizbox}[1][]{enhanced,breakable,colback=surface,colframe=rulec,
  boxrule=0.4pt,arc=2pt,left=9pt,right=9pt,top=8pt,bottom=8pt,#1}
\newtcolorbox{callbox}[2][accent]{enhanced,breakable,colback=#2,colframe=#2,
  boxrule=0pt,arc=0pt,borderline west={2.2pt}{0pt}{#1},
  left=9pt,right=9pt,top=7pt,bottom=7pt}
\newenvironment{callA}{\begin{callbox}[accent]{accentsoft}}{\end{callbox}}
\newenvironment{callW}{\begin{callbox}[warnc]{warnsoft}}{\end{callbox}}
\newenvironment{callS}{\begin{callbox}[stopc]{stopsoft}}{\end{callbox}}
\newcommand{\ct}[1]{{\mono\fontsize{7.2}{8}\selectfont\color{ink2}\MakeUppercase{#1}}\par\vspace{2pt}}

% a layer card: count, heading, prose, family tags
\newcommand{\layer}[4]{%
\begin{card}
  \begin{minipage}[t]{13mm}
    \raggedright{\mono\fontsize{17}{18}\selectfont\color{accent}#1}\\[-1pt]
    {\mono\fontsize{6}{7}\selectfont\color{ink3}FAMILIES}
  \end{minipage}\hfill
  \begin{minipage}[t]{\dimexpr\linewidth-16mm\relax}
    {\sffamily\bfseries\fontsize{10.5}{12.5}\selectfont #2}\par\vspace{2.5pt}
    {\color{ink2}\fontsize{9}{11.6}\selectfont #3\par}\vspace{4pt}
    \raggedright #4
  \end{minipage}
\end{card}\vspace{3pt}}

% a question and answer card
\newcommand{\qa}[2]{%
\begin{card}
  {\sffamily\bfseries\fontsize{10.3}{12.5}\selectfont #1}\par\vspace{3pt}
  {\color{ink2}\fontsize{9}{11.6}\selectfont #2\par}
\end{card}\vspace{3pt}}

% a flow step: tag, name, text
\newcommand{\stepbox}[3]{%
  \eyebrow{#1}\par\vspace{1pt}
  {\sffamily\bfseries\fontsize{9.6}{11}\selectfont #2}\par\vspace{1.5pt}
  {\color{ink2}\fontsize{8.4}{10.4}\selectfont #3\par}}

% ---------------------------------------------------------------- lists
\newenvironment{tight}
  {\begin{itemize}[leftmargin=13pt,itemsep=4pt,topsep=4pt,parsep=0pt,
     label=\textcolor{ink3}{\footnotesize$\bullet$}]\color{ink2}}
  {\end{itemize}}

% a legend row: chip in a fixed left column, meaning beside it
\newcommand{\legrow}[2]{%
  \noindent\begin{minipage}[t]{21mm}\raggedright #1\end{minipage}\hspace{2mm}%
  \begin{minipage}[t]{\dimexpr\linewidth-23mm\relax}
    {\fontsize{9}{11.4}\selectfont\color{ink2}#2\par}
  \end{minipage}\par\vspace{4pt}}

% --------------------------------------------------------------- tables
\newcolumntype{N}{>{\mono}r}
\newcolumntype{L}{>{\mono}l}
\newcommand{\thd}[1]{{\mono\fontsize{6.8}{8}\selectfont\color{ink3}\MakeUppercase{#1}}}
\newenvironment{sheettable}[1]
  {\par\vspace{5pt}{\mono\fontsize{7.4}{9}\selectfont\color{ink3}\raggedright #1\par}\vspace{2.5pt}%
   \begingroup\arrayrulecolor{rulec}\scriptsize
   \setlength{\tabcolsep}{3pt}\renewcommand{\arraystretch}{1.22}}
  {\endgroup\par}
\newcommand{\tnote}[1]{\par\vspace{4pt}{\fontsize{8.4}{10.6}\selectfont\color{ink3}#1\par}}

% -------------------------------------------------------- bars and meters
% Speedup bars. Track is 104 mm wide and spans 0 to 100x, so a value v
% in x units is v*1.04 mm. Change both numbers together to rescale.
\newcommand{\barline}[3]{%
  \begin{tikzpicture}[x=1mm,y=1mm,baseline=0.2mm]
    \fill[surf2] (0,0) rectangle (104,2.9);
    \fill[#2] (0,0) rectangle (#1*1.04,2.9);
  \end{tikzpicture}\hspace{2mm}{\mono\fontsize{7}{8}\selectfont\color{ink2}#3}}
\newcommand{\speedrow}[5]{%
  \noindent\begin{minipage}[t]{23mm}\raggedright
    {\mono\fontsize{7.4}{9}\selectfont\color{ink2}#1}\end{minipage}%
  \begin{minipage}[t]{\dimexpr\linewidth-23mm\relax}
    \barline{#2}{stwo}{#3}\\[1.3pt]
    \barline{#4}{sone}{#5}
  \end{minipage}\par\vspace{3.5pt}}
\newcommand{\baraxis}{%
  \noindent\hspace{23mm}\begin{tikzpicture}[x=1mm,y=1mm]
    \draw[rulesoft,line width=0.4pt] (0,2) -- (104,2);
    \foreach \p/\l in {0/0x,26/25x,52/50x,78/75x,104/100x}
      {\node[anchor=north,inner sep=1pt,text=ink3] at (\p,1.6)
        {\mono\fontsize{6.6}{7}\selectfont \l};}
  \end{tikzpicture}\par\vspace{2pt}}

% Prefill against decode meters. Track is 164 mm, so a share s in
% percent is s*1.64 mm.
\newcommand{\meterrow}[5]{%
  \noindent{\mono\fontsize{7.6}{9}\selectfont #1}\hfill
  {\mono\fontsize{7}{8}\selectfont\color{ink3}#2}\par\vspace{1.5pt}
  \noindent\begin{tikzpicture}[x=1mm,y=1mm]
    \fill[surf2] (0,0) rectangle (164,5.4);
    \fill[sone] (0,0) rectangle (#3*1.64,5.4);
    \fill[stwo] (#3*1.64+0.5,0) rectangle (164,5.4);
    \node[anchor=west,inner sep=0pt,text=white] at (2,2.7)
      {\mono\fontsize{6.8}{8}\selectfont #4};
    \node[anchor=east,inner sep=0pt,text=white] at (162,2.7)
      {\mono\fontsize{6.8}{8}\selectfont #5};
  \end{tikzpicture}\par\vspace{5pt}}
\newcommand{\swatch}[1]{\tikz[baseline=-0.4ex]\fill[#1] (0,0) rectangle (2.4mm,2.4mm);}
\newcommand{\vlegend}[1]{{\fontsize{8.6}{10}\selectfont\color{ink2}#1\par}\vspace{5pt}}
\newcommand{\viztitle}[2]{%
  {\sffamily\bfseries\fontsize{10.3}{12}\selectfont #1}\par
  {\fontsize{8.4}{10}\selectfont\color{ink3}#2\par}\vspace{7pt}}

\begin{document}
\sloppy

% ===================================================== MASTHEAD
\begingroup
\eyebrow{Hardware and software co-design \textperiodcentered{} program reference}\par\vspace{5pt}
{\rmfamily\bfseries\fontsize{27}{30}\selectfont CHEETAH Program Sheet\par}\vspace{6pt}
{\fontsize{10.6}{14}\selectfont\color{ink2}
One joint linear program picks the accelerator and its whole schedule at
once: MAC count, bandwidth, memory technology, on-chip capacity,
per-operation mapping, residency, rematerialization and the KV window are
decided together, and every design ships with a certificate that brackets
its own latency.\par}\vspace{7pt}
{\mono\fontsize{7.8}{10}\selectfont\color{ink3}
\textbf{Twelve workloads} designed end to end \quad
\textbf{68} row families \quad
\textbf{28k to 664k} variables \quad
\textbf{0.35M to 5.4M} nonzeros \quad
\textbf{HiGHS/CORGIsolver}\par}
\vspace{5pt}
{\color{ink}\hrule height 1.4pt}
\endgroup
\vspace{10pt}

% ===================================================== LEGEND
\begin{card}
{\mono\fontsize{7.4}{9}\selectfont\color{ink3}PROVENANCE OF EVERY NUMBER ON THIS SHEET\par}\vspace{4pt}
\legrow{\chipM}{Built and counted on \sheetdata{}, on the RTX 5090 host \idt{kevin}, from the deployed Stage-1 code that produced the reports. The build-only probe reproduced each report's rows and columns and each certificate to every printed digit.}
\legrow{\chipR}{Read out of the twelve \idt{*_e2e.json} run reports of 2026-09-21, which survive on \idt{kevin} and \idt{mazeppa}.}
\legrow{\chipN}{Recorded in \idt{DESIGN_deepseek.md} or the project log, not re-derived here.}
\legrow{\chipL}{Produced on the two A100 pods and the H100 pod, which no longer answer. The outputs were never mirrored to a surviving host.}
\end{card}

% ===================================================== THE PROGRAM
\section{The program, layer by layer \chipM}
\lede{Every row of the master program carries a family name, so the layers
below are the real tags from the source, not a retelling. Counted from
\idt{joint_lp.py}: 69 family declarations, 68 distinct names, \idt{op_mass}
opened twice.}

\layer{11}{Chip identity and selectors}{%
One chip for the whole workload. Exactly one PE type, one memory technology
and one bandwidth point are chosen, and every operation's mapping must be
compatible with that same chip. The on-chip global memory is tied to the
installed tile slices by an equality, and the bandwidth and memory technology
are coupled so an illegal pair cannot be assembled.}{%
\rowfam{pe_type_onehot}\rowfam{memtech_onehot}\rowfam{bw_onehot}\rowfam{bw_x_memtech}\rowfam{cgm_is_tile_slices}\rowfam{mode_matches_chip}\rowfam{mode_matches_chip_opt}\rowfam{bw_matches_chip}\rowfam{bw_matches_chip_opt}\rowfam{vector_lanes_le_chip}\rowfam{vector_lanes_domain}}

\layer{6}{Physical gates}{%
The six pre-RTL manufacturability limits sit inside the program as rows, not
as a filter applied afterwards, so an illegal chip is infeasible rather than
discarded later. Reticle area is an equality against the area model, which is
deliberate: a free area epigraph would let the power-density limit relax
itself.}{%
\rowfam{gate_reticle_area}\rowfam{gate_sram_macro}\rowfam{gate_crossbar}\rowfam{gate_crossbar_tiles}\rowfam{gate_power_density}\rowfam{gate_total_power}}

\layer{10}{Mapping selection and mass conservation}{%
Each operation picks one implementation mode out of the cached library: a PE
aspect, a global-buffer slice size, a buffer-keep policy and a double-buffering
depth. The mass rows conserve work across alternative implementations of the
same operation, and the copy rows stop the chip from running more tile copies
than it installs.}{%
\rowfam{op_mass}\rowfam{block_onehot}\rowfam{converter_onehot}\rowfam{mass_le_block}\rowfam{retain_le_block}\rowfam{tile_copies_le_installed}\rowfam{group_tile_copies_le_installed}\rowfam{copy_domain_tiles}\rowfam{copy_needs_chip_feature}\rowfam{mode_x_bw}}

\layer{6}{On-chip capacity}{%
What is resident at a moment has to fit. Retained tensors are charged at full
size, the KV window and the staging buffers are charged separately, and a
tensor's release interval has to cover every later consumer, so a long skip
cannot be kept for free.}{%
\rowfam{residency_capacity}\rowfam{kv_window_capacity}\rowfam{capacity_workspace}\rowfam{staging_capacity}\rowfam{fast_capacity}\rowfam{release_interval_covers_later_consumers}}

\layer{13}{Residency credits, the McCormick layer}{%
This is where traffic saved by keeping a tensor on chip gets credited, and it
is the only place a software decision multiplies another software decision.
The credit is bounded by the chosen mode, by what was actually retained, and
by the KV window, in both the per-operation and the grouped form. The products
are relaxed by McCormick envelopes, which is also where the relaxation gap
comes from.}{%
\rowfam{residency_credit_le_mode}\rowfam{residency_credit_le_retained_add}\rowfam{residency_credit_le_retained_out}\rowfam{residency_credit_group_add}\rowfam{residency_credit_group_out}\rowfam{kv_credit_le_mode}\rowfam{kv_credit_le_window}\rowfam{kv_credit_group_le_window}\rowfam{fast_credit_le_mode}\rowfam{fast_credit_le_class}\rowfam{fast_kv_credit_le_mode}\rowfam{fast_kv_credit_le_held}\rowfam{fast_kv_credit_le_held_group}}

\layer{4}{KV format and tiering}{%
The KV cache's numeric format is a decision with its own one-hot, and the
on-chip window can never exceed what the chosen format allows. The fast-tier
rows couple the same window to a faster memory class when one is installed.}{%
\rowfam{kv_format_onehot}\rowfam{kv_window_le_format}\rowfam{fast_share_le_format}\rowfam{fast_kv_le_format}}

\layer{3}{Rematerialization}{%
An activation can be recomputed instead of stored. A tensor cannot be both
retained and recomputed, a recomputed input must be read from DRAM, and the
recompute time is charged against the operation that pays for it.}{%
\rowfam{remat_plus_retained_le_1}\rowfam{remat_inputs_in_dram}\rowfam{remat_time_ge_share}}

\layer{12}{Time, the epigraph layer}{%
No row says what an operation's time is. Each row says the time variable is at
least something: at least the compute term, at least the memory term, at least
the conversion and vector work, at least what the host link and the service
limit allow. Minimizing the sum then pushes each time variable down onto the
binding term, which is how the maximum of compute and memory stays linear.}{%
\rowfam{op_time}\rowfam{op_time_vector}\rowfam{mem_time}\rowfam{conv_time}\rowfam{conv_hull}\rowfam{conv_le_point}\rowfam{load_path}\rowfam{host_link_phase}\rowfam{little_latency_credited}\rowfam{overlap_by_buffering}\rowfam{persp_vector}\rowfam{service_limit}}

\layer{3}{Quality and admission}{%
A family that changes numerics can only be used inside a stated accuracy
budget, drawn from an accuracy-vetted menu in the contract. Stage 1 pins
precision at the baseline, so these rows are present and inert.}{%
\rowfam{quality_family_use}\rowfam{quality_budget}\rowfam{quality_joint_exempt}}

\begin{callA}
\ct{What makes the hardware continuous}
Installed MACs and bandwidth are not picked off a grid of allowed chips. The
reciprocal cost $D/x$ of a mode with fixed demand $D$ running at rate $x$ is
handled by perspective tangent planes, so the rate stays a continuous variable
inside a pure linear program.
\vspace{4pt}
\[
\begin{array}{@{}r@{\;}c@{\;}l@{\qquad}l@{}}
  \tau &\ge& \dfrac{D z^{2}}{x_k}
    & \text{\footnotesize the convex hull of the disjunction}\\[9pt]
  \tau &\ge& D\!\left(\dfrac{2z}{g}-\dfrac{x_k}{g^{2}}\right)
    & \text{\footnotesize the tangent plane at anchor } g\\[9pt]
  \mathrm{slack} &=& \dfrac{D\,(x_k-gz)^{2}}{g^{2}x_k}\;\ge\;0
    & \text{\footnotesize so no plane ever cuts off a feasible point}
\end{array}
\]
\vspace{2pt}
Anchor spacing comes from an error target rather than taste: a geometric grid
with ratio $(1+\sqrt{\varepsilon})/(1-\sqrt{\varepsilon})$ meets the target
everywhere, which over a 256\tx{} range is 29 anchors at
$\varepsilon=10^{-2}$ and 89 at $\varepsilon=10^{-3}$.
\end{callA}

% ===================================================== SIZES
\section{Program sizes and nonzeros \chipM}
\lede{Built fresh on \sheetdata{}, one program per workload, with no solve.
Columns and rows match each run report; nonzeros and the variable-group count
were never recorded before, so these are new numbers.}

\begin{sheettable}{Stage-1 root program per workload. Build seconds are this probe; solve seconds are from the 2026-09-21 reports.}
\begin{adjustbox}{max width=\linewidth}
\begin{tabular}{L N N N N N N N N N N N}
\toprule
\thd{Workload} & \thd{Ops} & \thd{Variables} & \thd{Var groups} & \thd{Rows}
 & \thd{Equality} & \thd{Inequality} & \thd{Nonzeros} & \thd{nnz/row}
 & \thd{Density} & \thd{Build s} & \thd{Solve s}\\
\midrule
llama1\_13b      &  30 &  28,284 &   229 &   103,575 &  2,235 &   101,340 &   352,121 & 3.40 & 0.0120\% & 0.4 &     3.0\\
llama1\_65b      &  30 &  28,284 &   229 &   103,575 &  2,235 &   101,340 &   352,121 & 3.40 & 0.0120\% & 0.4 &     4.4\\
llama3\_70b      &  35 &  32,189 &   264 &   120,030 &  2,605 &   117,425 &   407,576 & 3.40 & 0.0105\% & 0.5 &     5.4\\
deepseek\_v3     & 136 & 111,070 &   971 &   452,421 & 10,079 &   442,342 & 1,527,767 & 3.38 & 0.0030\% & 1.3 &    18.7\\
efficientnet\_b1 & 116 & 212,201 & 1,074 &   503,594 &  8,599 &   494,995 & 1,773,095 & 3.52 & 0.0017\% & 1.1 &   394.3\\
efficientnet\_b2 & 116 & 212,201 & 1,074 &   503,594 &  8,599 &   494,995 & 1,773,095 & 3.52 & 0.0017\% & 1.2 &   255.4\\
efficientnet\_b3 & 131 & 239,771 & 1,212 &   568,847 &  9,709 &   559,138 & 2,002,904 & 3.52 & 0.0015\% & 1.3 &   548.3\\
efficientnet\_b4 & 161 & 294,911 & 1,488 &   699,353 & 11,929 &   687,424 & 2,462,522 & 3.52 & 0.0012\% & 1.5 &   586.4\\
efficientnet\_b5 & 195 & 357,499 & 1,801 &   847,356 & 14,445 &   832,911 & 2,983,807 & 3.52 & 0.0010\% & 1.6 & 1,320.1\\
efficientnet\_b6 & 225 & 412,639 & 2,077 &   977,862 & 16,665 &   961,197 & 3,443,425 & 3.52 & 0.0009\% & 1.8 & 1,469.1\\
efficientnet\_b7 & 274 & 502,797 & 2,528 & 1,191,118 & 20,291 & 1,170,827 & 4,194,519 & 3.52 & 0.0007\% & 2.1 & 2,506.5\\
sd15\_unet       & 346 & 664,221 & 3,251 & 1,533,481 & 25,619 & 1,507,862 & 5,416,291 & 3.53 & 0.0005\% & 2.6 & 2,427.0\\
\bottomrule
\end{tabular}
\end{adjustbox}
\end{sheettable}
\tnote{Three things are worth reading off this table. The programs are
extremely sparse, about 3.4 to 3.5 nonzeros per row at every scale, which is
what a program built from one-hots, epigraphs and two-term McCormick envelopes
looks like. Rows outnumber variables roughly three to one, because most of the
content is bounding rows rather than new decisions. And solve time does not
track size: EfficientNet-B1 is a quarter of SD-UNet's size yet takes 394
seconds against 2,427, because what costs time is the branch tree and the
degeneracy, not the row count.}
\tnote{On reproduction: the four LLM programs now carry one more inequality row
than their reports, because the SRAM-macro gate row was added at 19:00 on
2026-09-21, after those four runs. EfficientNet and SD-UNet match exactly.
Every certificate value reproduced to every printed digit, for example
Llama-1 13B at $L=14047.098659307783$ ms and $U=14055.489119933172$ ms.}

% ===================================================== INPUTS
\section{Timeloop and Vizier no longer in the Loop \chipM}
\lede{No mapper and no black-box search
runs inside the optimization.}

\begin{tcbraster}[raster columns=2,raster column skip=5pt,raster equal height,
  raster before skip=4pt,raster after skip=4pt]
\begin{tcolorbox}[cardstyle,breakable=false]
\stepbox{offline, once}{Timeloop}{Native mapper at search depth 1000 builds an
implementation library per workload. 40,922 records are cached on this host
across the twelve libraries.}
\end{tcolorbox}
\begin{tcolorbox}[cardstyle,breakable=false]
\stepbox{offline, once}{Contract}{Task shape, model config, numerics, accuracy
menu and routing statistics, with source hashes.}
\end{tcolorbox}
\end{tcbraster}

\begin{tcbraster}[raster columns=3,raster column skip=5pt,raster equal height,
  raster before skip=0pt,raster after skip=6pt]
\begin{tcolorbox}[cardstyle,breakable=false,colframe=accent,
  borderline west={2pt}{0pt}{accent}]
\stepbox{in the loop}{Joint LP}{Reads cached records as fixed coefficients.
Chooses chip and schedule together.}
\end{tcolorbox}
\begin{tcolorbox}[cardstyle,breakable=false,colframe=accent,
  borderline west={2pt}{0pt}{accent}]
\stepbox{in the loop}{Branch tree}{Bandwidth, then PE type by memory
technology, each node its own LP.}
\end{tcolorbox}
\begin{tcolorbox}[cardstyle,breakable=false,colframe=accent,
  borderline west={2pt}{0pt}{accent}]
\stepbox{in the loop}{Recovery and gates}{Legal array multiples and bank
geometry, six gates, exact lift, re-solve at the fixed chip.}
\end{tcolorbox}
\end{tcbraster}

\qa{\vno{not in the loop}Timeloop}{%
Timeloop is referenced only by the library builder, the replay checker and the
schedule evaluator. The four modules that run during a solve, \idt{run_e2e.py},
\idt{joint_lp.py}, \idt{lp_eval.py} and \idt{sched_eval.py}, mention it only in
comments explaining which Timeloop dataspace a coefficient came from. A library
is built once per workload, then every solve reads it. Library records on this
host: DeepSeek 5,196, EfficientNet 3,640 each, Llama 2,215 to 2,458, SD-UNet
3,358.}

\qa{\vno{not in the loop}Vizier}{%
There is no Vizier import anywhere in the solve path. The only occurrences in
the design package are in the baseline builder, which calls
\idt{vizier_full.scale_tl_to_bw} as the published baseline's evaluator, not as
a search. Vizier did drive the older published pipeline, where it sampled chips
and each trial was evaluated separately, and it was last used that way for the
SD-UNet experiment of 2026-09-20. In the current design the LP sizes the chip,
which is the whole point: putting a sampler around it would move hardware back
outside the program.}

\qa{\vpart{designed, not yet wired}The outer loop}{%
The planned role for a proposer is narrow: read the solved program's duals and
reduced costs, propose new \emph{implementations} to add to the library, and
let the LP price them. It never picks hardware, batch size or budgets. Nothing
of this has run yet, and the stopping rule is deliberately not ``the proposer
found nothing''.}

% ===================================================== MODELS
\section{Models that have been run \chipR}
\lede{Twelve workloads, all twelve designed end to end with the frozen Stage-1
program, all passing six of six physical gates. Uniform INT8 W8A8 with an
8-bit KV cache on both the design and the baseline, so no quantization gain
hides inside a co-design number.}

\begin{sheettable}{Task definition and the recovered chip, per workload. TPU-v3 reference for scale: 65,536 MACs, 900 GB/s, 32 MiB on-chip, \mmsq{132.2} in this area model.}
\begin{adjustbox}{max width=\linewidth}
\begin{tabular}{L l N N N l N l N N}
\toprule
\thd{Workload} & \thd{Task} & \thd{Tiles} & \thd{MACs} & \thd{vs TPU-v3}
 & \thd{PE type} & \thd{GB/s} & \thd{Memory} & \thd{On-chip MiB} & \thd{Gates}\\
\midrule
deepseek\_v3     & B 32, prompt 4096, 512 tokens & 738 & 188,928 & 2.88\tx & \mono sq16x16, 256 KiB     & 3,200 & \mono HBM  & 200.5 & 6/6\\
llama1\_13b      & B 32, prompt 1536, 512 tokens & 789 & 201,984 & 3.08\tx & \mono sq16x16, 256 KiB     & 3,200 & \mono HBM  & 199.2 & 6/6\\
llama1\_65b      & B 32, prompt 1536, 512 tokens & 789 & 201,984 & 3.08\tx & \mono sq16x16, 256 KiB     & 3,200 & \mono HBM  & 199.2 & 6/6\\
llama3\_70b      & B 32, prompt 4096, 512 tokens & 742 & 189,952 & 2.90\tx & \mono tcgen05\_32x8, 256 KiB & 3,200 & \mono HBM  & 200.5 & 6/6\\
sd15\_unet       & 50 denoising steps, UNet batch 2 & 877 & 224,512 & 3.43\tx & \mono sq16x16, 256 KiB   & 1,800 & \mono GDDR & 244.1 & 6/6\\
efficientnet\_b1 & one forward pass, batch 1     & 941 & 240,896 & 3.68\tx & \mono sq16x16, 256 KiB     & 2,400 & \mono HBM  & 236.4 & 6/6\\
efficientnet\_b2 & one forward pass, batch 1     & 940 & 240,640 & 3.67\tx & \mono wgmma8x32, 256 KiB   & 2,400 & \mono HBM  & 236.3 & 6/6\\
efficientnet\_b3 & one forward pass, batch 1     & 938 & 240,128 & 3.66\tx & \mono wgmma8x32, 256 KiB   & 2,400 & \mono HBM  & 236.3 & 6/6\\
efficientnet\_b4 & one forward pass, batch 1     & 956 & 244,736 & 3.73\tx & \mono wgmma8x32, 256 KiB   & 1,800 & \mono HBM  & 244.0 & 6/6\\
efficientnet\_b5 & one forward pass, batch 1     & 953 & 243,968 & 3.72\tx & \mono wgmma8x32, 256 KiB   & 1,800 & \mono HBM  & 243.9 & 6/6\\
efficientnet\_b6 & one forward pass, batch 1     & 928 & 237,568 & 3.62\tx & \mono wgmma8x32, 256 KiB   & 1,800 & \mono GDDR & 244.0 & 6/6\\
efficientnet\_b7 & one forward pass, batch 1     & 931 & 238,336 & 3.64\tx & \mono wgmma8x32, 256 KiB   & 1,800 & \mono GDDR & 243.9 & 6/6\\
\bottomrule
\end{tabular}
\end{adjustbox}
\end{sheettable}
\tnote{The LP is not told which PE aspect to prefer. It chose square 16x16
tiles for the LLMs and SD-UNet, the tall 8x32 aspect for EfficientNet-B2
upward, and the 32x8 aspect for Llama-3 70B, where that choice was certified to
matter: the same program relaxed to 124.9 seconds with \idt{tcgen05_32x8}
against 135.9 with \idt{wgmma8x32} and 145.2 with \idt{sq16x16}.}

\subsection{After Stage 1}
\begin{tight}
\item \textbf{\color{ink}Stage 2} \chipN{} added the quantization and sparsity
families on three workloads: Llama-1 13B gained 1.37\tx{} over its Stage-1
design with W4 plus 4-bit KV, Llama-1 65B 1.24\tx, and Llama-3 70B only
1.014\tx{} because W4 is not admitted for it.
\item \textbf{\color{ink}The v5 master program} is merged and gated, 18 of 18
gate checks passed, and its campaign ran on the two A100 pods. \chipL{} Those
outputs lived in \idt{e2e/v5/} on the pods, no surviving host holds a copy, and
the pods now refuse connections. The code, the contracts and the libraries
survive, so the campaign is repeatable.
\item \textbf{\color{ink}}
\end{tight}

% ===================================================== PERFORMANCE
% \section{Performance \chipR}
% \lede{Two baselines are in play, and the gap between them is the single most
% important caveat on this page. Quote the one you mean.}

% \begin{callW}
% \ct{Read this before quoting a speedup}
% \textbf{Baseline A} is the published recipe. Its Timeloop runs used a mapping
% constraint that places work on only one axis of the PE array, built from an
% 8x8 reference rescaled as if it had 256 MACs, so its compute term is about 32
% times the full-array ideal. It inflates every speedup by roughly seven to ten
% times.
% \par\vspace{5pt}
% \end{callW}

% \needspace{80mm}
% \begin{vizbox}[breakable=false]
% \viztitle{Speedup over TPU-v3, by baseline convention}{Same axis, same unit,
% two baselines. One linear scale from 0 to 100\tx.}
% \vlegend{\swatch{stwo}\ vs baseline A, the published recipe \qquad
%          \swatch{sone}\ vs the same-evaluator TPU-v3-like chip}
% \speedrow{deepseek\_v3}{30.60}{30.60\tx}{4.20}{4.20\tx, +320\%}
% \speedrow{llama1\_13b}{37.74}{37.74\tx}{3.75}{3.75\tx, +275\%}
% \speedrow{llama1\_65b}{35.29}{35.29\tx}{3.63}{3.63\tx, +263\%}
% \speedrow{llama3\_70b}{44.96}{44.96\tx}{4.76}{4.76\tx, +376\%}
% \speedrow{sd15\_unet}{96.78}{96.78\tx}{9.49}{9.49\tx, +849\%}
% \speedrow{efficientnet\_b7}{65.85}{65.85\tx}{10.44}{10.44\tx, +944\%}
% \baraxis
% \tnote{Baseline A figures are from the twelve reports. Same-evaluator figures
% are from the notes of 2026-09-22; the EfficientNet family spans 6.13\tx{} at
% B1 to 10.44\tx{} at B7.}
% \end{vizbox}

\subsection{Prefill against decode}
{\color{ink2}The task is one prefill pass over the whole prompt followed by the
generation of 512 tokens, and the split between the two is not similar across
models. Llama-3 70B is 85\% prefill. DeepSeek-V3 is 59\% decode.\par}
\vspace{6pt}

\needspace{50mm}
\begin{vizbox}[breakable=false]
\viztitle{Where the task's milliseconds go}{Share of total latency, on the
recovered design. Values are milliseconds.}
\vlegend{\swatch{sone}\ prefill \qquad
         \swatch{stwo}\ decode, 512 tokens in four context bins}
\meterrow{deepseek\_v3}{136,204 ms total}{40.86}{55,653 \textperiodcentered{} 40.9\%}{80,551 \textperiodcentered{} 59.1\%}
\meterrow{llama1\_13b}{14,055 ms total}{56.84}{7,988 \textperiodcentered{} 56.8\%}{6,067 \textperiodcentered{} 43.2\%}
\meterrow{llama1\_65b}{61,081 ms total}{58.47}{35,716 \textperiodcentered{} 58.5\%}{25,365 \textperiodcentered{} 41.5\%}
\meterrow{llama3\_70b}{126,337 ms total}{85.02}{107,412 \textperiodcentered{} 85.0\%}{18,925 \textperiodcentered{} 15.0\%}
\end{vizbox}

\begin{sheettable}{Per-phase latency, utilization and the speedup each phase earns against baseline A.}
\begin{adjustbox}{max width=\linewidth}
\begin{tabular}{L N N N N N N N N}
\toprule
\thd{Workload} & \thd{Prefill ms} & \thd{Decode ms} & \thd{Prefill MFU}
 & \thd{Decode MFU} & \thd{Prefill x} & \thd{Decode x} & \thd{Total x}
 & \thd{Tokens/s}\\
\midrule
deepseek\_v3 &  55,652.8 & 80,551.5 & 0.617 & 0.083 & 62.87\tx &  8.31\tx & 30.60\tx &   120.3 (re-run)\\
llama1\_13b  &   7,988.5 &  6,067.0 & 0.406 & 0.181 & 52.47\tx & 18.35\tx & 37.74\tx & 1,165.7\\
llama1\_65b  &  35,715.8 & 25,364.7 & 0.455 & 0.215 & 45.11\tx & 21.46\tx & 35.29\tx &   268.2\\
llama3\_70b  & 107,412.3 & 18,925.0 & 0.474 & 0.342 & 47.17\tx & 32.39\tx & 44.96\tx &   129.7\\
sd15\_unet   & \multicolumn{2}{c}{\mono one denoising phase, 230.6} & \multicolumn{2}{c}{\mono 0.776} & \multicolumn{2}{c}{\mono n/a} & 96.78\tx & n/a\\
\multicolumn{1}{@{}l}{\mono efficientnet\_b1 to b7} & \multicolumn{2}{c}{\mono single pass, 0.0129 to 0.3336} & \multicolumn{2}{c}{\mono 0.221 to 0.475} & \multicolumn{2}{c}{\mono n/a} & \multicolumn{1}{r}{\mono 34.32\tx{} to 65.85\tx} & n/a\\
\bottomrule
\end{tabular}
\end{adjustbox}
\end{sheettable}

\begin{callA}
\ct{The finding the phase split carries}
Prefill reaches a utilization of 0.41 to 0.62 and earns 45\tx{} to 63\tx{}
against baseline A. Decode reaches 0.083 to 0.342 and earns 8\tx{} to 32\tx.
Decode is bound by bytes streamed per token, not by arithmetic, and on
DeepSeek-V3 that traffic is dominated by mixture-of-experts weight reload
rather than by the KV cache, whose latent attention keeps its cache small. The
same program says so in the utilization column: 0.083 on DeepSeek decode is
the memory wall, measured.
\end{callA}
% ===================================================== KV
\section{KV cache active across timesteps\chipR}
\lede{Yes. It is modelled per decode step, it grows with context inside the
task, and its on-chip window is a decision the program makes rather than a
parameter it is handed.}

\qa{\vyes{active}Four context bins per task, carried through the objective}{%
Generation of 512 tokens is modelled in four context bins, so the attention
cost is charged against a growing history rather than one average length. The
reports carry a separate latency for each bin, and attention compute is charged
over the history, not over one token.}

\qa{\vyes{a decision}The on-chip KV window}{%
Capacity rows \idt{kv_window_capacity} and \idt{kv_window_le_format} make the
resident fraction of the cache a variable, bounded by the chosen numeric
format, competing with weights and activations for the same on-chip bytes. The
window is time-multiplexed by layer, so one layer's worth of positions is
resident at a time.}

\qa{\vpart{conventions to quote}Lossless, with no prefix credit}{%
The KV budget is pinned to the full history, so no eviction policy is assumed.
The prefix-cache reuse credit is switched off, after it was found to cancel
essentially all decode KV traffic from a 0.6\% hit rate. Free eviction, which
had discarded positions regardless of context, was removed as well.}

\begin{sheettable}{Decode latency per context bin, and the fraction of the KV cache the design keeps on chip in each bin.}
\begin{adjustbox}{max width=\linewidth}
\begin{tabular}{L l N N N N N N N}
\toprule
\thd{Workload} & \thd{Attention} & \thd{Bin 0 ms} & \thd{Bin 1 ms}
 & \thd{Bin 2 ms} & \thd{Bin 3 ms} & \thd{Growth}
 & \thd{KV on chip, bin 0} & \thd{bin 3}\\
\midrule
llama1\_13b  & MHA, 400 KiB/token  &  1,415.9 &  1,483.1 &  1,550.4 &  1,617.6 & +14.2\% & 9.79e-5 & 7.89e-5\\
llama1\_65b  & MHA                 &  6,018.4 &  6,233.6 &  6,448.8 &  6,663.9 & +10.7\% & 3.06e-5 & 2.47e-5\\
llama3\_70b  & GQA, 160 KiB/token  &  4,580.3 &  4,829.7 &  4,615.6 &  4,899.4 &  +7.0\% & 7.23e-4 & 6.62e-4\\
deepseek\_v3 & MLA, latent cache   & 19,960.2 & 20,020.9 & 20,403.8 & 20,166.6 &  +1.0\% & 3.59e-3 & 3.28e-3\\
\bottomrule
\end{tabular}
\end{adjustbox}
\end{sheettable}
\tnote{The growth column is the honest signature of a live KV cache:
multi-head attention at 400 KiB per token makes Llama-1 13B's last bin 14.2\%
slower than its first, while DeepSeek's latent attention grows 1.0\% because
its cache is tiny and its decode cost sits in expert weight reload instead.
Llama-3 70B is not monotone across bins, which is the mode mix changing between
bins rather than noise. The resident fractions are small on every workload, so
the program is choosing to spend its on-chip bytes on weights and activations
rather than on the cache, and saying so with a number.}

% ===================================================== SOLVE
\section{How it is solved \chipM}
\lede{One LP solver, called many times, inside a shallow branch and bound whose
branches are the discrete choices of the chip. There is no mixed-integer solver
anywhere in the pipeline.}

\subsection{Branch and bound}
{\color{ink2}Branch and bound searches a discrete space without enumerating it.
Relax the discrete choices so the problem becomes a linear program and solve
that: the answer is a \emph{lower bound} on the true optimum, because the
relaxation can do anything the real problem can and more. If the relaxation
happens to come out discrete, it is the answer and the search is over.
Otherwise pick an undecided choice and \emph{branch}: make one child problem
per value, each a new LP with that value pinned. Any complete design found
along the way gives an \emph{upper bound}, the incumbent. A child whose own
lower bound is already worse than the incumbent can be discarded without being
explored, since nothing inside it can beat what is already in hand, and that
pruning is what keeps the tree small. The search ends when the best remaining
lower bound is close enough to the incumbent, and the pair of numbers is the
certificate.\par}

\subsection{This program's tree}
\begin{tight}
\item \textbf{\color{ink}Root.} Solve the full relaxation. If the bandwidth
selector comes out integral and the gap is within the tolerance, stop. This
happened on all four LLM workloads, where the reports record ``root bandwidth
integral''.
\item \textbf{\color{ink}First level, bandwidth.} Eight delivered-bandwidth
points from 225 to 3,200 GB/s, one pinned LP each. Pruning is visible in the
reports: on EfficientNet-B1 the 225, 450 and 675 GB/s branches were cut with
``pruned: bound 0.0379892 $\ge$ U''.
\item \textbf{\color{ink}Second level, PE type by memory technology.} Eighteen
nodes, three PE aspects by two slice sizes by three memory technologies,
entered only when the gap is still open. Llama-3 70B and SD-UNet both needed
it, and SD-UNet pruned ten of its eighteen nodes on bounds.
\item \textbf{\color{ink}Sub-branching.} In the v5 program a solved branch with
an open gap is split again on its most fractional remaining selector, KV
format, block, feature or converter, two levels deep.
\item \textbf{\color{ink}Upper bound.} Every candidate is turned into a real
chip: legal array multiples and SRAM bank geometry, the six gates checked, the
reciprocals and products set to their true values and every original row
re-checked, then the fixed-chip scheduling LP re-solved. Incumbents are carried
across branches so the upper bound never gets worse.
\item \textbf{\color{ink}Lower bound.} The safe dual bound of the relaxation,
with duals sign-clipped on one-sided rows, or the minimum over branches when
the enumeration is complete.
\end{tight}

\begin{sheettable}{Certificates as reported. The gap is what the design is guaranteed to be within.}
\begin{tabular}{L N N N l}
\toprule
\thd{Workload} & \thd{Lower bound} & \thd{Upper bound} & \thd{Gap}
 & \thd{Branching needed}\\
\midrule
llama1\_13b      & 14,047.10 ms & 14,055.49 ms & 0.06\% & none, root integral\\
llama1\_65b      & 61,005.97 ms & 61,080.51 ms & 0.12\% & none, root integral\\
deepseek\_v3     & 135,919.79 ms & 136,204.25 ms & 0.21\% & none, root integral\\
llama3\_70b      & 124,883.99 ms & 126,337.33 ms & 1.15\% & PE type by memory technology, 18 nodes\\
efficientnet\_b1 & 0.012820 ms & 0.012927 ms & 0.83\% & bandwidth, 3 of 8 pruned\\
efficientnet\_b7 & 0.332239 ms & 0.333645 ms & 0.42\% & bandwidth, 8 nodes\\
sd15\_unet       & 223.26 ms & 230.58 ms & 3.17\% & bandwidth, then 18 nodes, 10 pruned\\
\bottomrule
\end{tabular}
\end{sheettable}
\tnote{EfficientNet B2 to B6 land between 0.25\% and 0.54\%. SD-UNet's 3.17\%
is the loosest in the set and is mode rounding rather than chip choice.}

\subsection{Every relaxation technique in the program}
\begin{tcbraster}[raster columns=2,raster column skip=5pt,raster equal height=rows,
  raster before skip=5pt,raster after skip=4pt]
\begin{tcolorbox}[cardstyle,breakable=false]
{\sffamily\bfseries\fontsize{10}{12}\selectfont McCormick envelopes}\par\vspace{3pt}
{\color{ink2}\fontsize{8.8}{11}\selectfont Used only where a decision
multiplies a decision, which here is a residency credit times a mode or
bandwidth selector. For $z=xy$ with $x\in[l,u]$ and $y\in[p,q]$ the four rows
are the standard envelope, two lower and two upper. This is the dominant source
of the remaining gap.\par}
\end{tcolorbox}
\begin{tcolorbox}[cardstyle,breakable=false]
{\sffamily\bfseries\fontsize{10}{12}\selectfont Perspective tangent cuts}\par\vspace{3pt}
{\color{ink2}\fontsize{8.8}{11}\selectfont Reciprocal hardware cost, time as
demand over rate, outer-approximated by tangent planes at geometric anchors.
Each plane is valid, so the bound is never wrong, only loose, and adding
anchors tightens it. This is what keeps MACs and bandwidth continuous without a
grid of allowed chips.\par}
\end{tcolorbox}
\begin{tcolorbox}[cardstyle,breakable=false]
{\sffamily\bfseries\fontsize{10}{12}\selectfont Epigraph variables}\par\vspace{3pt}
{\color{ink2}\fontsize{8.8}{11}\selectfont Twelve families of them. Every
maximum in the cost model, compute against memory, host link against local
service, is a set of lower-bounding rows on one time variable rather than a
nonlinear max.\par}
\end{tcolorbox}
\begin{tcolorbox}[cardstyle,breakable=false]
{\sffamily\bfseries\fontsize{10}{12}\selectfont One-hot selectors with disaggregated products}\par\vspace{3pt}
{\color{ink2}\fontsize{8.8}{11}\selectfont PE type, memory technology,
bandwidth, block, KV format and converter are one-hot. Selector-times-continuous
products are disaggregated exactly, per branch, rather than written as a single
bytes-over-bandwidth product, which is where an earlier version leaked a large
spurious saving.\par}
\end{tcolorbox}
\begin{tcolorbox}[cardstyle,breakable=false]
{\sffamily\bfseries\fontsize{10}{12}\selectfont Chord rows, for the inner bound}\par\vspace{3pt}
{\color{ink2}\fontsize{8.8}{11}\selectfont The tangent planes bound from
outside. Chords bound from inside, and are only valid once every selector is
pinned, so they belong to the achievable side of the sandwich rather than to
the certified lower bound.\par}
\end{tcolorbox}
\begin{tcolorbox}[cardstyle,breakable=false]
{\sffamily\bfseries\fontsize{10}{12}\selectfont Exact lift}\par\vspace{3pt}
{\color{ink2}\fontsize{8.8}{11}\selectfont The step that makes the upper bound
trustworthy: reciprocals and products are set to their true values and every
original row is re-checked. Residuals land around $10^{-16}$ when it passes. An
LP that converged is not evidence that its products were exact.\par}
\end{tcolorbox}
\end{tcbraster}

% ===================================================== BORROWED
\section{Shape Aggregation \chipN}
\lede{The following is new and semi-experimental.}

\begin{tight}
\item \textbf{\color{ink}Shape aggregation.} We cluster operation instances and into unique operation
cluster shapes: the shape-invariant part of the program, tiling, roofline,
envelopes and hardware, is built once per unique shape and weighted by
multiplicity, while the time-indexed residency stays per instance. On
DeepSeek-V3 that collapses 903 operations to 13 unique shapes, which is what
makes the program faster to solve. But may limit search space. Further analysis needed.
\item \textbf{\color{ink}Dual-guided rounding.} The row duals guide which fractional selector to round
and which branch to open first.
\item \textbf{\color{ink}The shadow-price outer loop.} Traditionally the rule for reading
prices off a solved relaxation and acting on them becomes the template for the
planned implementation designer, which reads duals and proposes library columns.
\end{tight}
\tnote{}

% ===================================================== OPEN
\section{Update Fri Oct 9}
\section{Large Scale Joint Inference Solve}
\lede{One joint program sizes hardware and schedules software together, so the
optimizer exploits interactions that separate sizing, kernel mapping, batching and
memory planners miss. Choices span compute/SRAM/bandwidth, per-operation kernel
mappings, optional batching and fusion, KV placement, transfers, eviction and
recomputation. We index execution and state over ordered event intervals to model time.}

\begin{tight}
\item \textbf{\color{ink}First accepted generated solve.} The reduced causal
DeepSeek LP certifies at 168.26 s with a 0.0179\% interval, inside the 0.1\% bar,
over all 344,158 variables and 687,307 rows, with independent replay. CORGI
supplied the hardware values, recovery a feasible primal, and a small CPU
certificate LP the strong lower bound; the raw CORGI iterate stayed infeasible.
This validates a hybrid generated solver, not raw GPU convergence.
\item \textbf{\color{ink}Why it is hard.} Representation: $\sim$506M logical
invocations make full expansion impractical, and streaming coefficients and operator native methods alone does
not remove the vectors. Search: optional batches create an enormous subset universe
to prune or certify away without enumerating. Correctness: individually reasonable
scheduling, memory and hardware decisions can be jointly infeasible against the
readiness, storage, transfer, capacity and service constraints.
\item \textbf{\color{ink}Approach.} Keep the full target implicit, optimize a
smaller generated working program, and certify what stays outside it. Independent
decisions stay independent; exact state aliases compress bookkeeping without forcing
same-shaped invocations to share decisions.
\item \textbf{\color{ink}Mechanisms proved.} Persistent configurations;
sparse-exception pricing; membership-adjustable certificates; temporal certificates
with sparse state --- together producing the accepted interval on the complete,
unchanged reduced causal LP.
\item \textbf{\color{ink}Where we are.} A checked, coupled reduced-task solve, not
isolated examples, but full DeepSeek is still unsolved. The full-input registry now
accounts for every family.
\item \textbf{\color{ink}Next step.} Extend the proved machinery to the full schema:
host/DRAM/SRAM copies, variable-vector hardware, the larger execution library,
complete omitted-alternative certificates, and a replayable full-task primal
witness. 
\item \textbf{\color{ink}End result.} A compact replayable solution over the entire
declared workload with independently checked bounds within 0.1\%; a recovered chip
and schedule with physical validity and its own relaxation gap; matched TPUv3
reference, CHEETAH on fixed TPUv3, and joint co-design, split into prefill, decode
and total; plus runtime, convergence diagnostics, feature coverage and durable
checkpoints.
\end{tight}


% ===================================================== TRAINING ROADMAP
\clearpage
\section{CHEETAH-compiler: DeepSeek training on $N$ B200s \chipN}
\lede{A plan for a fixed cluster. The hardware is given.
The compiler jointly chooses the complete repeated training schedule and emits
an executable distributed plan. The items below are targets, not measured B200 results. The model is DeepSeek-V3 671B, not a smaller distilled model. The architecture configuration has 61 transformer layers, hidden dimension 7,168, and 256 routed experts with eight selected per token.}

\begin{callA}
\ct{Research thesis}
Parallelism, expert placement, Blackwell kernels, activation retention,
communication and microbatch scheduling compete for the same resources.
Optimize them together over a complete training step, then check the emitted
execution against the declared model and numerical contract. Builds from the fixed Cost plan (calibrate twice works).
\end{callA}

\subsection{Freeze the training contract}
\begin{tight}
\item Fix the exact DeepSeek configuration, optimizer and state precision,
sequence length, non-padding global token batch, accumulation convention and
allowed numerical error. Declare $N$, GPU memory available to the job, the
NVLink domains and the inter-node network topology.
\item Preserve the learned router, routed token assignments, token weights,
optimizer-update boundaries and stochastic semantics in the first experiment.
Expert replication must aggregate replica gradients and synchronize weights;
it is an execution choice, not independent expert training.
\item Compare complete synchronized training steps, including communication
and optimizer updates. Verify loss, gradients and updated parameters against
the reference within the stated tolerance, then check short-run training quality.
\end{tight}

\subsection{One complete, repeated training graph}
\begin{card}
\ct{Task components}
Attention forward and backward; router and top-$k$; token permutation;
dispatch; expert GEMM1, activation and GEMM2; combine; shared experts;
activation-gradient tasks $B_X$; weight-gradient tasks $B_W$; gradient
reduction/scatter; optimizer updates; retained-state transfers and recomputation.
Separate $B_X$ from $B_W$ so weight-gradient work can move where dependencies allow.
\end{card}
The primary objective is full-step time at a fixed global batch:
\[
 \min T_{\mathrm{step}},\qquad
 \mathrm{throughput}=\frac{\text{non-padding tokens per step}}{T_{\mathrm{step}}}.
\]
Within a legal accumulation window, a repeated microbatch schedule may minimize
its initiation interval $II$. Charge fill, drain, reductions and the optimizer
boundary explicitly: a low $II$ alone does not establish a faster training step.
No forward task may read updated parameters before its declared update boundary.

\subsection{Seven training decision families}
\layer{01}{Module parallelism and topology}{Choose compatible DP, TP, EP, PP
and CP layouts from a finite validated menu. Attention, shared experts and routed
experts may use different layouts; charge every layout conversion and ensure all
process groups fit the same physical cluster.}{\rowfam{parallelism}\rowfam{layout_conversion}\rowfam{topology}}
\layer{02}{Communication as scheduled work}{Split dispatch, expert compute and
combine into dependent chunks. Choose chunk size, transfer path, communication
SM allocation, in-flight depth and buffering. Compare NCCL, DeepEP and custom
transfer implementations only where their training semantics are validated.
Model NVLink and inter-node links separately.}{\rowfam{dispatch}\rowfam{combine}\rowfam{link_capacity}}
\layer{03}{Blackwell implementation columns}{Build measured forward, $dX$ and
$dW$ expert GEMMs; MLA forward/backward; normalization/residual; router;
permutation; gradient reduction and optimizer columns. Expose tiles, layouts,
warp specialization, TMA pipelines, supported \idt{tcgen05}/TMEM usage,
CTA clusters and fusion. Every column must have an emitter and a correctness check.}{\rowfam{implementation}\rowfam{tile}\rowfam{resource}}
\layer{04}{Full training memory}{Jointly select retention, legal reload,
recomputation, fused recomputation and optional offload. Include parameters,
master weights, Adam moments, gradients, activations, dispatch buffers,
communication scratch, attention state, RNG state and kernel workspace.
HBM is the backing store on this target; ``spill'' must name its source tier.}{\rowfam{residency}\rowfam{rematerialization}\rowfam{capacity}}
\layer{05}{Microbatch pipeline}{With the global batch fixed, choose microbatch
size/count and legal overlap of forward, $B_X$, $B_W$ and communication.
Search double/triple buffering and eventually modulo schedules, with explicit
buffer lifetimes and update barriers.}{\rowfam{microbatch}\rowfam{precedence}\rowfam{buffer_lifetime}}
\layer{06}{Expert placement and replication}{Place experts across the fixed
cluster; optionally replicate hot experts inside a fast communication domain.
Charge replica parameters, optimizer/gradient ownership, synchronization traffic
and straggler load. Select among replicas without changing token-to-expert semantics.}{\rowfam{expert_placement}\rowfam{replication}\rowfam{load}}
\layer{07}{Global resource schedule}{Choose task start times and implementation
combinations under measured shared-resource envelopes. Communication and compute
can overlap only when SM, HBM, link, scratch and dependency constraints permit;
concurrent latency is profiled rather than inferred from isolated peaks.}{\rowfam{schedule}\rowfam{overlap}\rowfam{contention}}

\clearpage
\subsection{Program structure and solver roles}
\lede{Use three timescales to keep the search tractable. Joint optimization is
the goal, and staged search is our initial implementation strategy.}
\qa{Structural search}{Choose compatible parallelism, placement,
replication, region partition, microbatch configuration and column menus.
Use LP relaxations with branching/refinement for discrete choices. OCS solves
continuous relaxations only.}
\qa{Repeated-step planning}{Given structural choices, solve the structured
schedule and memory relaxation through operator products. Recover a discrete
feasible schedule, check all original constraints and replay it. A restricted
subproblem bound does not certify the unrestricted global search.}
\qa{Cheap runtime adaptation}{Use a compiled portfolio or bounded local policy
for route-count fluctuations, replica assignment and queue admission. Prices
may guide choices, but correctness and capacity guards remain explicit. Include
fallback cost to avoid a large optimization solve on every training iteration.}

\begin{card}
\ct{Capacity and shared-resource constraints}
For every GPU $g$ and event time $t$,
\[
 M_g^{\mathrm{params}}+M_g^{\mathrm{optimizer}}+M_g^{\mathrm{grad}}
 +M_g^{\mathrm{live}}(t)+M_g^{\mathrm{comm}}(t)
 +M_g^{\mathrm{workspace}}(t)\le H_g.
\]
Include temporary allocations and replicated state once, under explicit
ownership rules. Concurrent communication and compute allocations satisfy
$n_g^{\mathrm{comm}}+n_g^{\mathrm{compute}}\le n_g^{\mathrm{SM}}$;
SMEM, TMEM, registers and warp limits are checked at their hardware scopes.
Link constraints are per link/domain, with collective synchronization and
transfer dependencies represented in the graph.
\end{card}

\subsection{What OCS can contribute}
\begin{tight}
\item \textbf{\color{ink}Operator-native planning.} Preserve forward/adjoint
access through program construction, normalization, solving and reporting.
Independently check the recovered plan and any claimed objective bound.
\item \textbf{\color{ink}Batched portfolios and warm starts.} Solve related
sequence lengths, microbatch menus, route-load scenarios and memory budgets.
Use observed histograms and stress cases; marginal load percentiles alone need
not describe a valid joint routing scenario.
\item \textbf{\color{ink}Prices and sensitivities.} Read duals to identify
valuable reductions in link traffic, workspace or kernel time on the fixed
hardware. Treat sensitivities as local to the continuous model; discrete plan
changes and degeneracy can invalidate a naive derivative.
\item \textbf{\color{ink}Column generation.} Use reduced cost
$\mathrm{RC}_j=c_j-\lambda^\top a_j$ under the declared sign convention to
prioritize candidate kernels. Generate, verify, benchmark and insert a column,
then re-solve. A LM proposer remains outside the trusted checker.
\end{tight}

\subsection{A hierarchical Blackwell executor}
\begin{callA}
\ct{Execution boundary}
A global plan drives attention, router/dispatch, expert forward, combine,
attention backward, expert $dX$, expert $dW$, gradient communication and optimizer
regions. Use CUDA graphs or a validated device controller between regions;
use persistent kernels within regions where they improve measured performance.
\end{callA}
% Within eligible regions, use TMA producers, tensor-core consumers, TMEM,
% explicit barriers, buffering and supported CTA clustering. Across GPUs, emit
% explicit communication tasks. Establish forward progress, dependency visibility
% and resource residency before enabling cross-region overlap; a whole-cluster
% single megakernel is not a prerequisite.

\subsection{The overlap pattern to search}
Allow the current microbatch's expert GEMMs to overlap a later microbatch's
dispatch, an earlier microbatch's combine, and ready $dW$/gradient reductions.
Attention from another ready layer can fill slack. This pattern is a candidate,
not a policy imposed in advance: the optimizer must respect data readiness,
collective ordering, contention, memory lifetimes and the parameter-version barrier.

\clearpage
\subsection{Bounds, prediction and measured execution}
For a declared implementation menu and compatible numerical formats, construct
valid lower bounds from unavoidable work, per-resource traffic and precedence:
\[
 L_{\mathrm{train}}=\max\left\{
 \frac{F_{\min}}{\sum_g P_g},\;
 \max_g\frac{B_{g,\min}^{\mathrm{HBM}}}{W_g^{\mathrm{HBM}}},\;
 \max_{\ell}\frac{B_{\ell,\min}}{W_{\ell}},\;
 T_{\mathrm{critical},\min},\;T_{\mathrm{imbalance},\min}
 \right\}.
\]
The minima must be justified across legal choices; traffic from one selected
plan is not automatically a global lower bound. Use actual topology and any
required cut constraints. Distinguish this analytic floor from the checked
relaxation bound $L_{\mathrm{model}}$, the feasible modeled plan time
$U_{\mathrm{plan}}$, and measured time $T_{\mathrm{measured}}$.
\begin{sheettable}{Three diagnostics; all interpreted within their declared execution and cost model.}
\begin{tabular}{p{40mm}p{120mm}}
\toprule
\thd{Diagnostic} & \thd{Interpretation}\\
\midrule
$T_{\mathrm{baseline}}/L_{\mathrm{train}}$ & Available modeled headroom, if the floor is valid.\\
$U_{\mathrm{plan}}/L_{\mathrm{model}}$ & Certified planning gap when both cover the same search space.\\
$T_{\mathrm{measured}}/U_{\mathrm{plan}}$ & Realization error; check calibration, contention and runtime overhead.\\
\bottomrule
\end{tabular}
\end{sheettable}

\subsection{Development gates}
\begin{tight}
\item \textbf{\color{ink}1. Contract and baseline.} Freeze the target and build
a tuned production baseline from the selected Megatron/DeepSpeed/DeepEP stack.
Measure synchronized full steps, route histograms, memory peaks and stage times.
\item \textbf{\color{ink}2. B200 calibration.} Profile the actual cluster;
measure isolated and concurrent column costs. Do not reuse H100 coefficients
without validation. Include the inter-node network when $N$ exceeds one domain.
\item \textbf{\color{ink}3. Kernel and memory closure.} Qualify expert F/$dX$/$dW$,
MLA F/B, dispatch/combine and optimizer paths. Implement full persistent-state
accounting. Re-test joint rematerialization in a memory-binding training regime;
the manuscript's earlier gain is evidence to reproduce, not a B200 prediction.
\item \textbf{\color{ink}4. One-domain schedule.} Add microbatch and chunked
communication scheduling; recover and replay feasible plans. Verify numerical
semantics, capacity and deadlock-free repeated execution.
\item \textbf{\color{ink}5. Cluster expansion.} Add expert placement/replication,
compatible module-specific layouts and their conversions. Validate scaling on
the actual $N$-B200 topology and stressed route loads.
\item \textbf{\color{ink}6. Joint search and evaluation.} Integrate structural
search, step planning and dual-guided column generation. Report compile/setup
cost, amortization, modeled gaps and measured execution together.
\item \textbf{\color{ink}7. Differentiable routing, separately.} Only after the
semantics-preserving compiler works, study a plan-aware router with
$\mathcal L(\theta,\ell)+\mu T^*(\ell)$. State the differentiability assumptions,
relaxation and gradient estimator; evaluate model quality and time to target
loss because this experiment changes training behavior.
\end{tight}

\subsection{Flagship experiment and ablations}
\begin{card}
\ct{Same cluster, same declared training contract}
Compare a tuned production baseline with Blackwell columns, then add joint
rematerialization, communication scheduling, placement/replication,
module-specific layouts and the global schedule. Pair the waterfall with
leave-one-out ablations and a matched staged planner using the same column
library, so gains from kernels and gains from joint decisions can be distinguished.
\end{card}
Report non-padding tokens/s, synchronized step time, explicitly defined MFU,
HBM/link bytes, exposed communication, idle SM time, peak memory, optimization
time and amortization, checked modeled gap and realization error. Record run
variance, routing stress behavior and numerical checks. For any router-changing
extension, also report convergence and time to matched quality.

\begin{callW}
\ct{Status and scope}
This section records a proposed fixed-B200 training direction informed by
\emph{Compiler\_megakernel\_ICML2027}, especially its training graph,
rematerialization and hardware-explicit execution columns. It asserts no new
B200 measurements, no universal optimality and no completed solver capability.
The INT8 inference conventions in the preceding program sheet do not apply to
this training plan; its precision is part of the contract to be frozen.
\end{callW}

% ===================================================== FOOTER
\vspace{14pt}
{\color{ink}\hrule height 1.4pt}
\vspace{7pt}
{\mono\fontsize{7.6}{10}\selectfont\color{ink3}
Sources. Program structure and sizes: \idt{systems/cheetah_joint/joint_lp.py}
and a build-only probe run on kevin, \sheetdata. Performance, chips and
certificates: the twelve \idt{*_e2e.json} reports of 2026-09-21 on kevin.
Numerics throughout are INT8 W8A8 with an 8-bit KV cache on both design and
baseline. Latency is in milliseconds at the design's own clock. Speedups are
ratios of task latency under one baseline convention; percentages are the
same ratios less one.\par\vspace{4pt}
.
Data measured \sheetdata, source last edited \sheetrev.\par}

\end{document}
